% =============================================================================
% CHAPTER 6: Reflection on Learning
% =============================================================================
% SPDX-FileCopyrightText: 2026 Frank Langbein <frank@langbein.org>
% SPDX-License-Identifier: LPPL-1.3c
%
% Suggested sections: Critical narrative of the experience; Analysis and application;
% Implications for future practice. Focus on transferable learning and (where
% applicable) ``double loop learning'': impact on assumptions and concepts, not
% only skills. Optional alternatives: ``What I would do differently''; ``Transferable
% skills and concepts''. Omit this chapter if not required by your institution;
% comment out the \chapter and \input in main.tex.
% =============================================================================

This chapter reflects on what was learned from the project: how assumptions and decisions were re-evaluated and what would be done differently in future.

% Talk about this:
% Applying double-loop learning to my methodology, my initial assumption was that an algorithm's slowness was always caused by a massive Branch-and-Bound search tree. 
% However, tracking the node count metric dismantled this assumption: the solver was getting stuck on the root node simply trying to process the sheer volume of dense constraints.
% I had to do some learning and understand that matrix sparsity (as used in the Ship model) is just as critical to performance as search-space pruning.

\section{Critical narrative of the experience}\label{sec:challenges}
% Key challenges, turning points, and how they affected your understanding.


\section{Analysis and application}\label{sec:reflection}
% How you analysed the experience and applied learning to decisions or concepts.


\section{Implications for future practice}\label{sec:implications}
% What you will take forward: habits, assumptions, or approaches in future work.

